\documentclass[10pt,a4paper]{beamer}
\usepackage[utf8]{inputenc}
\usepackage{wrapfig}

\usetheme{Copenhagen}

\title{FindAndTravel}
\author{Gianluca Lutero \and Klajdi Kola}
\institute{Universita' degli studi di Parma}
\date{}


\begin{document}

\begin{frame}
\maketitle
\end{frame}

\begin{frame}{Indice}
\begin{figure}

\begin{minipage}{0.47\textwidth}
   \begin{itemize}
      \item Obiettivi dell'applicazione
      \item Componenti Android utilizzate
      \item Funzionalita' dell'applicazione
      \item Sviluppi futuri
   \end{itemize}
\end{minipage}
\hfill 
\begin{minipage}{0.35\textwidth}

   \includegraphics[scale=1]{logo2.png}

\end{minipage}

\end{figure}
\end{frame}

\begin{frame}
\frametitle{Obiettivi dell'applicazione} 
    \begin{block}{Obiettivo}
    
    L'applicazione FindAndTravel si pone come obiettivo quello di fornire informazioni turistiche per punti di interesse dell'utente. 
    \end{block}      
Cosa puo' fare l'utente: 
   \begin{itemize}
    \item Cercare i punti di interesse di una localita'
    \item Ottenere informazioni per quei luoghi, in particolare:
    \begin{itemize}
         \item Nome del luogo
         \item Posizione
         \item Tipo
         \item Feedback degli utenti
    \end{itemize}        
    \item Lasciare un feedback   
   \end{itemize}    
\end{frame}

\begin{frame}{Componenti Android utilizzate}
\begin{figure}

\begin{minipage}{0.47\textwidth}
   \begin{itemize}
      \item Google Map
      \item Google Places per Android e Web
      \item SQLite
      \item Volley
      \item Location service
   \end{itemize}
\end{minipage}
\hfill 
\begin{minipage}{0.35\textwidth}

   \includegraphics[scale=0.1]{Android_robot.png}

\end{minipage}

\end{figure}

\end{frame}

\begin{frame}{Funzionalita' di FindAndTravel}
\framesubtitle{Ciclo di vita di FindAndTravel}
  \begin{figure}
   \centering
   \includegraphics[scale=0.2]{CicloDiVitaApp.png}
  \end{figure}
\end{frame}

\begin{frame}{Funzionalita' di FindAndTravel}
\framesubtitle{Menu' principale}
\begin{figure}
\begin{minipage}{0.45\textwidth}
 \begin{block}{Menu'}
 
 Dal menu principale l'utente ha la possibilita' di accedere alle varie sezioni dell'app attraverso i corrispondenti \texttt{ImageButton}.    
 \end{block}
\end{minipage}
\hfill
\begin{minipage}{0.35\textwidth}
  \includegraphics[scale=0.1]{menu.png}
\end{minipage}
\end{figure}

\end{frame}

\begin{frame}{Funzionalita' di FindAndTravel}
\framesubtitle{Mappa}
\begin{figure}
\begin{minipage}{0.45\textwidth}
 \begin{block}{Mappa}
  Sulla mappa sono evidenziati i punti d'interesse attraverso i marker cliccabili.
 \end{block}
\end{minipage}
\hfill
\begin{minipage}{0.35\textwidth}
  \includegraphics[scale=0.1]{mappa2.png}
\end{minipage}
\end{figure}

\end{frame}

\begin{frame}{Funzionalita' di FindAndTravel}
\framesubtitle{Dettagli punto d'interesse}
\begin{figure}
\begin{minipage}{0.45\textwidth}
 \begin{block}{Dettagli}
  Sono visualizzate le informazioni relative al punto d'interesse. Attraverso il bottone in alto è possibile salvare tra i preferiti il luogo mentre con quello in basso si accede a \textbf{Google Maps} sulla posizione indicata con la possibilità di vedere piu' commenti degli utenti e lasciarne uno proprio. 
 \end{block}
\end{minipage}
\hfill
\begin{minipage}{0.35\textwidth}
  \includegraphics[scale=0.1]{dettagli.png}
\end{minipage}
\end{figure}
\end{frame}

\begin{frame}{Funzionalita' di FindAndTravel}
\framesubtitle{Preferiti}
\begin{figure}
\begin{minipage}{0.45\textwidth}
 \begin{block}{Preferiti}
 Sono salvate tutte le posizioni preferite dell'utente visualizzate attraverso una lista. Toccando una riga della lista si accede alla sezione dettagli relativa.
 \end{block}
\end{minipage}
\hfill
\begin{minipage}{0.35\textwidth}
  \includegraphics[scale=0.1]{preferiti.png}
\end{minipage}
\end{figure}
\end{frame}


\begin{frame}{Funzionalita' di FindAndTravel}
\framesubtitle{Ricerca}
\begin{figure}
\begin{minipage}{0.45\textwidth}
 \begin{block}{Ricerca}
 In questa sezione è possibile ricercare i luoghi che non vengono visualizzati attraverso la mappa. Una volta selezionato il risultato desiderato si verrà riportati nella relativa sezione dettagli.
 \end{block}
\end{minipage}
\hfill
\begin{minipage}{0.35\textwidth}
  \includegraphics[scale=0.1]{ricerca.png}
\end{minipage}
\end{figure}
\end{frame}

\begin{frame}{Sviluppi futuri}
\begin{itemize}
\item Porting iOS
\item Maggiore integrazione con i servizi di Google
\item Miglioramento dell'interfaccia grafica
\end{itemize}
\end{frame}

\end{document}